% !TEX root = ../3.GSM.tex
\section{Theory}

\subsection{Capacity and Performance of the Mill}

The efficiency calculations of the mill use input particle size $D_{p1}$ (ft) and the size of 80\% of product after milling $D_{p2}$ (ft).

$P$ is capacity for milling material from a very large size to $D_p$ (volume/time, with $i = \infty$):
\begin{equation}
    P = K_b \sqrt{\frac{1}{D_p}}
\end{equation}

According to the definition, $W_i$ is the necessary energy to mill from a very large size down to 100\,\textmu m (kWh/ton). The relation between $W_i$ and $K_b$ (the Bond coefficient, which depends on the type of material and milling machine) is:
\begin{equation}
    60W_i = K_b \sqrt{\frac{1}{100 \times 10^{-3}}}
    \quad \Rightarrow \quad
    K_b = \frac{60W_i}{\sqrt{10}} \approx 19W_i
    \quad \Rightarrow \quad
    P = 19W_i \frac{1}{\sqrt{D_p}}
\end{equation}

where $P_1 = 19W_i / \sqrt{D_{p1}}$ and $P_2 = 19W_i / \sqrt{D_{p2}}$.

Milling capacity for crushing 1 ton of material over 1 minute from $D_{p1}$ to $D_{p2}$:
\begin{equation}
    P = P_2 - P_1 = 19W_i \left(\frac{1}{\sqrt{D_{p2}}} - \frac{1}{\sqrt{D_{p1}}}\right)
\end{equation}

$T$ is productivity (ton/min). Capacity for milling $T$ tons/min from $D_{p1}$ to $D_{p2}$:
\begin{equation}
    P = 19W_i \left(\frac{1}{\sqrt{D_{p2}}} - \frac{1}{\sqrt{D_{p1}}}\right) T \quad \text{(kW)}
\end{equation}

where $D_{p1}$, $D_{p2}$ are the sizes of the material and product in mm. If $P$ is multiplied by $4/3$, the consumption capacity of the motor is:
\begin{equation}
    P = UI\cos\varphi
\end{equation}

Where: $U$ -- voltage (V); $I$ -- amperage (A); $\cos\varphi$ -- power factor.

Milling efficiency:
\begin{equation}
    H = \frac{P}{P'} \times 100\%
\end{equation}

\subsection{Size Distribution}

\begin{equation}
    \frac{-d\varPhi}{dD_p} = K D_p^b
\end{equation}

Where: $\varPhi$ -- cumulative mass fraction on $D_p$; $D_p$ -- particle size; $K$, $b$ -- coefficients describing particle distribution.

Integrating from $\varPhi = \varPhi_1$ to $\varPhi = \varPhi_2$ corresponding to $D_p = D_1$ and $D_p = D_2$:
\begin{equation}
    \varPhi_2 - \varPhi_1 = \frac{K}{b+1}\left(D_{p1}^{b+1} - D_{p2}^{b+1}\right)
\end{equation}

Between the $n$-th and $(n+1)$-th sieve, using standard sieves where $D_{pn-1}/D_{pn} = r = \text{const}$:
\begin{equation}
    \Delta\varPhi_n = \varPhi_n - \varPhi_{n-1} = \frac{-K}{b+1}\left(D_{pn}^{b+1} - D_{pn-1}^{b+1}\right)
\end{equation}

Replacing $D_{pn-1} = r \cdot D_{pn}$:
\begin{equation}
    \Delta\varPhi_n' = \frac{K(r^{b+1}-1)}{b+1} D_{pn}^{b+1} = K' D_{pn}^{b+1}
    \quad \text{with} \quad K' = \frac{K(r^{b+1}-1)}{b+1}
\end{equation}

or equivalently:
\begin{equation}
    \log\Delta\varPhi_n = (b+1)\log D_{pn} + \log K'
\end{equation}

$K'$ and $b$ are determined by plotting $\Delta\varPhi_n$ vs.\ $D_{pn}$ on a log--log graph.

\subsection{Sieve Performance}

\begin{equation}
    E = \frac{J}{Fa} \times 100
\end{equation}

Where: $F$ -- initial material mass input to the sieve (g); $J$ -- material mass passing through the sieve (g); $a$ -- particle fraction that can pass through the sieve (\%).

$F \cdot a$ is determined in the laboratory by sifting $F$ to find $J_1$, then sifting the remainder $F - J_1$ to find $J_2$, and so on. The total $J_1 + J_2 + J_3 + \ldots$ will asymptote to $F \cdot a$. Sieve performance is 100\% if $J_1 = F \cdot a$.

\subsection{Mix Equation}

When mixing amount $a$ of component A with amount $b$ of component B to create a homogeneous mixture, the composition in an ideal mixture is:
\begin{itemize}
    \item With A: $\displaystyle C_A = \frac{a}{a+b}$
    \item With B: $\displaystyle C_B = \frac{b}{a+b}$
\end{itemize}

These components would be equal in any volume of the mixture in an ideal case. In reality, since mixing time cannot be infinite, components A and B will differ across different volume elements. To evaluate homogeneity, we consider the mean standard deviation:
\begin{equation}
    s_A = \sqrt{\frac{\sum_{i=1}^{N}(C_A - C_{iA})^2}{N-1}}
    \qquad
    s_B = \sqrt{\frac{\sum_{i=1}^{N}(C_B - C_{iB})^2}{N-1}}
\end{equation}

The smaller $s_A$, $s_B$ are, the more ideal the mixture. To evaluate the homogeneity in a more practical way, the mixing index is:
\begin{equation}
    I_s = \frac{\sigma_e}{s}
\end{equation}

with $\sigma_e$ the standard deviation of an ideal mixture:
\begin{equation}
    \sigma_e = \sqrt{\frac{C_A C_B}{n}}
\end{equation}

Therefore:
\begin{equation}
    I_s = \sqrt{\frac{C_A C_B (N-1)}{n \cdot \sum_{i=1}^{N}(C_A - C_{iA})^2}}
\end{equation}

where $n$ is the number of particles in the case of loose material mixing.
